\section{Discussion and Conclusion} 

\subsection{Discussion}
The results presented in Chapter~\ref{sec:results} demonstrated that the
prototype supports portable identity, structured social actions,
pull-based replication, account discovery and feed publication. These
results provide evidence that the requirements identified in Chapter~\ref{sec:methodology}
can be realised through the Git-native protocol proposed in
Chapter~\ref{sec:system-design}

The following sections discuss the major design decisions of the
protocol, compare them with existing approaches, and examine the
trade-offs revealed through implementation and evaluation.

\subsection{Identity and Portability}

The protocol's identity model builds on the portable public-key identities
used by Nostr: a user's Ed25519 public key identifies them, while the Git
repository stores and distributes their data \cite{fiatjaf2026}. This differs
from ActivityPub, where identity is associated with a server domain
\cite{w3cActivityPub2018}, and from Octotown, which uses a GitHub username
\cite{octotown2026}. Social4Git showed that repositories can move between
providers, but it did not provide a portable cryptographic identity
\cite{social4git2026}.

The identity-creation and publication workflows presented in
Chapter~\ref{sec:results} show how this works in the prototype. The client
adds the repository URL to Alice's signed profile, and
Figure~\ref{fig:results-publish-user} shows that publication uses the same
public key. This separates the key that identifies the user from the
repository location recorded in the profile. The Results demonstrate this
separation during publication, although they do not show a repository
migration between hosts.

The protocol supplements cryptographic identities with human-readable
handles. This approach is conceptually similar to the handle-based layer
used by AT Protocol, where a user-friendly identifier is mapped to a
portable cryptographic identity. The directory therefore improves
discoverability without changing the underlying identity model
\cite{atproto2026}.

The repository-replacement scenario also shows why the protocol pairs keys
with handles. A public key can verify a signed profile, but users need a
readable handle to find an account. AT Protocol uses a similar handle layer
alongside its cryptographic identifiers \cite{atproto2026}; in this
prototype, the social directory maps each handle to its registered key and
repository URL. The scenario demonstrates that Follow checks this mapping
before accepting the profile, so identity continuity depends on the
directory entry as well as the signature.

\subsection{Replication and Distribution}

The implemented protocol adopts a pull-based replication model built on
Git's native synchronisation mechanisms. A client fetches from configured
remotes, then checks candidate action files locally before accepting them.
This means Git transfers repository data while the client decides which
actions are valid for its local store.

The literature review highlighted different trade-offs in the surveyed
distribution approaches. ActivityPub uses server-to-server delivery, with
servers forwarding activities to followers; this can place increasing load
on servers as their follower networks grow \cite{w3cActivityPub2018}. Nostr
uses relays to forward events, but its relays do not track client state and
provide weak ordering guarantees \cite{fiatjaf2026,stackernews2026}. Secure
Scuttlebutt removes servers through peer-to-peer gossip, although its
append-only feeds grow as messages accumulate
\cite{ssbConsortium2026}.

Gitweets, Social4Git and GitSocial already use Git repositories and fetch
operations to share social data \cite{gitweets2026,social4git2026,gitsocial2026}.
Git-as-Federation-Transport also uses Git, but its repositories are managed
at the server level rather than for each user's social data
\cite{gittransport2026}.
This implementation builds on that approach by storing each interaction as
a signed action file, separate from the Git commit that carries it. After
fetching a remote, the client scans its \texttt{social/actions/} directory,
checks candidate signatures and accepts valid action files into the local
store. Git provides the repository transfer; the client applies the
protocol's acceptance checks.

The Results show this process in Alice's direct replication from Bob and
in the multi-hop scenario where Alice receives Charlie's post through Bob
(Figures~\ref{fig:results-replicated-bob-post} and
\ref{fig:results-alice-charlie}). Alice receives the same signed action
Charlie created and verifies it using Charlie's public key. Bob's repository
passes the action on, but Bob does not become its author. This demonstrates
multi-hop propagation for the tested workflow while keeping authorship with
the user who signed the action. The evaluation does not measure how this
workflow performs with larger repositories.

\subsection{Discovery and Usability}

Discovery emerged as one of the most significant gaps identified during
the protocol survey. ActivityPub provides discovery through WebFinger and
server infrastructure \cite{w3cActivityPub2018}, while AT Protocol uses
DID and handle services \cite{atproto2026}. Nostr provides portable
identities, but its relay model does not provide the same handle-to-repository
lookup \cite{fiatjaf2026}. Most Git-based systems require repository URLs
to be exchanged manually before users can connect to an account
\cite{gitweets2026,social4git2026,gitsocial2026}.

To address this limitation, the implemented protocol introduces a public
social directory that maps handles to repository URLs and public keys.
Users can look up a handle in the directory to find the repository URL,
rather than needing to obtain it from the account owner. During Follow, the
client checks the URL and fetched profile's handle and public key against
the directory entry. Signatures verify that the profile and actions match
the listed public key; the directory links that key to the handle.

The discovery and follow workflows presented in
Chapter~\ref{sec:results} demonstrate this approach as Alice follows Bob
(Figures~\ref{fig:results-directory-site} and
\ref{fig:results-follow-bob}). Unlike Gitweets, Social4Git and GitSocial,
where users need to obtain a repository URL before connecting, this
directory lets them find the URL using a handle
\cite{gitweets2026,social4git2026,gitsocial2026}. The user then provides
that URL to the Follow command. The client checks that it matches the URL
registered for the handle, then checks that the fetched profile has the
same handle, public key and repository URL as the directory entry.

The implementation also adds static-site publication through GitHub
Pages. The generated websites provide browser-accessible feeds displaying
profile information, posts, replies and likes (Figure~\ref{fig:results-feed-bob}).
These pages present data generated from the protocol's Git repositories and
JSON actions; they do not change how those actions are created or replicated.
This gives users a browser-based way to view their feed alongside the
command-line protocol client.

An important observation from the implementation is that usability
improvements frequently introduce additional trust assumptions. The
social directory significantly improves discoverability, and
browser-based feeds improve accessibility, but neither feature is
strictly required for protocol operation. Social actions can be stored and
exchanged through Git repositories, with cryptographic identities used to
sign and verify them. Account discovery uses the shared directory to map
handles to repositories and public keys. This separates social-data
exchange from the directory used to find accounts.

Compared with earlier Git-based systems, the directory and published feeds
add browser-based account discovery and feed viewing. Gitweets exposes Git
operations directly to users \cite{gitweets2026}, while Social4Git and
GitSocial use protocol-specific commands and require repository URLs to
connect \cite{social4git2026,gitsocial2026}. The prototype therefore makes
finding accounts and viewing posts more accessible, while core actions such
as following and replication still use the command-line client.

Nevertheless, the prototype remains primarily command-line oriented.
Users still perform operations such as following, replication and
publication through explicit commands. While this provides a higher
level of abstraction than systems which expose native Git commands
directly, it does not yet offer the level of accessibility expected of
mainstream social-networking applications.

\subsection{Social Vocabulary and Data Representation}

The protocol represents social interactions through a structured
vocabulary consisting of Posts, Replies, Likes and Follows. This design
was informed by the survey of existing systems. ActivityPub models
interactions through explicit activity types \cite{w3cActivityPub2018},
while Nostr distinguishes actions through event kinds \cite{fiatjaf2026}.

Unlike Gitweets, which uses commit messages for posts, and GitSocial, which
stores its action metadata in Git commit messages, the implemented protocol
stores social actions as independent JSON objects
\cite{gitweets2026,gitsocial2026}. This allows authorship, timestamps,
signatures and references to be represented as fields rather than inferred
from Git history. The client can validate these structured objects, and the
format provides a basis for extending the action vocabulary.

The action creation and interaction workflows presented in
Chapter~\ref{sec:results} show the common structure across Posts, Replies,
Likes and Follows. For example, a Reply stores the referenced post ID and
target handle alongside its content, while a Like stores the target post ID
and handle. Feed construction uses these fields to connect each interaction
to its post and author. This keeps social relationships in the action data
rather than relying on provider-specific functionality or inferring them
from Git history.

The use of JSON also provides extensibility. The protocol currently
supports only four action types, but the common action structure allows
new social interactions to be introduced without altering the existing
representation. The recent discussion regarding Boost or Retweet actions
provides a useful example: such an action could be introduced as an
additional JSON object referencing an existing post while preserving the
existing protocol architecture.

This demonstrates one of the key benefits of structured social data.
Earlier Git-based systems often used commit messages or
provider-specific features to represent interactions, limiting
flexibility and interoperability. By contrast, the implemented protocol
separates social activity from repository history and stores it using a
consistent machine-readable format. The result is a social vocabulary
that is easier to validate, extend and implement across different
clients.

\subsection{Overall Assessment}

The protocol was designed to combine strengths identified across both
Git-based and non-Git-based decentralised social systems while
addressing several recurring limitations identified during the
literature review.

Nostr demonstrates the use of portable public-key identity
\cite{fiatjaf2026}, while AT Protocol maps human-readable handles to
cryptographic identifiers \cite{atproto2026}. ActivityPub and Nostr
provide examples of structured social actions and event types
\cite{w3cActivityPub2018,fiatjaf2026}. Gitweets, Social4Git and GitSocial
show how repositories can support social-data storage and pull-based
replication \cite{gitweets2026,social4git2026,gitsocial2026}. The prototype
brings these ideas together: users store and exchange social actions
through Git repositories, while account discovery uses a shared directory.

The results presented in Chapter~\ref{sec:results} indicate that this
approach is feasible in practice. The prototype successfully
demonstrated portable identity, structured JSON actions, pull-based
replication, discovery through human-readable handles, multi-hop
content propagation and browser-accessible static feeds. Collectively,
these capabilities address many of the limitations identified in the
survey of existing protocols and subsequently formalised as requirements
in Chapter~\ref{sec:methodology}.

However, the implementation also revealed that some challenges are
transformed rather than eliminated. Identity portability reduces
dependence on hosting providers but increases reliance on trusted
handle-to-key mappings. Pull-based replication reduces infrastructure
requirements but introduces replication latency. Browser-based feed
publication improves accessibility but leaves core protocol operations
dependent on a command-line client. These trade-offs suggest that
decentralisation should not be viewed as the removal of all
dependencies, but rather as a redistribution of trust and
responsibility within the system.

These observations provide context for the security considerations,
limitations and future directions discussed in the following sections.


\subsection{Security Considerations}

The preceding sections discuss how the prototype separates public-key
identity from repository location and verifies actions during replication.
This section considers what those checks demonstrate, which identity claims
they support, and where the prototype still relies on trust outside the
signatures themselves.

\subsubsection{Protected Threats}

The prototype uses Ed25519 signatures to check that profiles and actions
have not changed since they were signed. The private key is stored in the
client's local identity state, while the profile and actions contain the
public key and signature needed for verification. As described in
Chapter~\ref{sec:implementation}, the same signing and verification
mechanisms are used across these data types. \texttt{ProfileVerifier}
checks profile signatures and \texttt{ActionVerifier} checks action
signatures. The site publisher also verifies actions again when building
the feed before including them in the published data. A valid signature
shows that the signed data matches the corresponding public key. It does
not, on its
own, prove that the key belongs to a particular handle. The social
directory supplies that handle-to-key link.

The Results chapter tests action integrity by changing the content of a
signed post while leaving its signature unchanged. Replication rejects the
modified action, as shown in Figures~\ref{fig:results-tampered-action} and
\ref{fig:results-reject-tamper}. The impersonation scenario changes the
post's author field to Alice's public key but keeps Bob's signature. The
client rejects it because the signature does not verify against the claimed
author key (Figures~\ref{fig:results-impersonated-action} and
\ref{fig:results-reject-impersonation}). These results show how the client
checks both that action content has not been altered and that its claimed
author key matches the signing key.

The repository-replacement scenario examines whether someone could take
over a handle by creating a new repository after the original repository
becomes unavailable. A replacement profile can have a valid signature
because it was signed with the new owner's key. This only shows that the
new owner controls that key. It does not show that the key belongs to the
existing handle. The directory records the public key associated with
each handle, and Follow checks the profile against that record. In the
test, the replacement profile uses the same handle and repository URL but
a different key. Follow rejects it because the directory still lists the
original key (Figures~\ref{fig:results-directory-original-key},
\ref{fig:results-replacement-profile} and
\ref{fig:results-reject-replacement}). This demonstrates how the client
can detect an attempted takeover when the new repository presents a
different key for an existing handle. The test models this conflict using
a replacement profile. It does not close and recreate a live hosted
repository.

The multi-hop result shows a different protection. When Alice receives
Charlie's action through Bob, she verifies it with Charlie's key. Bob can
pass the action on, but the action remains attributed to Charlie
(Figure~\ref{fig:results-alice-charlie}). These examples show the checks
working in the tested situations, rather than proving that every possible
attack is prevented.

These signatures add an authorship check that the surveyed Git-based
systems Social4Git and GitSocial do not provide. Git's commit history can
show that repository content changed, but without action signatures it
does not establish which account owner created a particular post
\cite{social4git2026,gitsocial2026}. The prototype therefore uses Git for
storage and replication while its own signature checks link each action
to the key named as its author.

\subsubsection{Trust Assumptions}

Signature verification does not require a server to approve each post or
replication step. The prototype's account discovery has a separate trust
requirement, however. The social directory is currently administered by
the project owner and records which public key and repository URL belong
to each handle. Follow relies on that record when it checks a profile. If
the directory entry is changed or its administrator registers the wrong
key, a valid signature alone cannot correct the handle-to-key mapping.
The repository-replacement result demonstrates that the client rejects a
key that conflicts with the registered entry, assuming that entry is
trustworthy. The client's publishing code also refuses to update an
existing handle with a different registered key.

The prototype also uses the \texttt{canonical-social} repository to provide
starter repository files and the site template used during publication,
as described in Chapter~\ref{sec:implementation}. This creates a separate
trust in the source of those files. A change to the site template could
affect how a published feed is rendered, but it does not change whether
profile and action signatures validate. The directory and template
repositories are therefore supporting services around the protocol, with
different roles and consequences if either is changed.

A signed identity document at a well-known URI is one possible alternative
for publishing a handle-to-key association. It would change where users
obtain that association, but would also make the URI's domain and its
ownership history part of the trust decision. This is a possible direction
for future work rather than a mechanism evaluated in this prototype.

\subsubsection{Security Limitations}\label{sec:limitation-security}

The current protocol provides integrity and key-based authorship verification through
digital signatures, but it does not encrypt social actions. Posts, replies, likes 
and follows are stored as readable JSON documents within repositories, meaning that 
anyone with access to those repositories can view their contents. 
For public social-media use cases, authenticity and integrity are often sufficient because
users primarily need to verify who created a post and whether it has been modified.

However, these guarantees are not sufficient for private communication. Private posts, 
direct messages or restricted-access social groups require confidentiality in addition 
to authenticity and integrity. Encryption protects content from being read by unauthorised 
users, whereas digital signatures only verify authorship and detect modification.

The sAT protocol provides a useful point of comparison because it demonstrates that confidentiality
can be achieved in a decentralised, static-hosted environment by encrypting content and 
distributing decryption keys to authorised followers \cite{sat2026}. However, the literature review 
also highlighted the additional complexity associated with this approach, 
including key distribution, key rotation and decryption overhead \cite{sat2026}. 
The present prototype therefore prioritises simplicity of replication and verification, 
at the cost of not supporting confidential social content.

Nevertheless, the sAT approach demonstrates that encryption is feasible and may be useful 
as an optional feature for private content rather than being applied to all social actions.
One possible direction is the use of shared session or group keys for authorised participants,
reducing the need to distribute a separate decryption key for every post. 
This possibility is discussed further in the Future Work section.

The repository-replacement check also depends on the directory maintaining the correct public key
for each registered handle. The current prototype assumes that the public key associated with
a handle does not change after registration. As a result, if a legitimate user loses their private
key and generates a new key pair, the new public key would not match the key currently stored in 
the directory. The client would therefore reject the updated profile in the same way that it 
rejects a repository-replacement attack, because it cannot distinguish between a genuine key change
and an attacker presenting a different key for an existing handle. The repository-replacement results
(Figures~\ref{fig:results-replacement-profile} and
\ref{fig:results-reject-replacement}) therefore highlight a trade-off between identity continuity
and identity recovery. The current
design prioritises continuity by rejecting unexpected key changes, but a practical deployment 
would also require a controlled mechanism for key rotation or account recovery.

The protocol also has no built-in moderation, blocking, spam filtering or Sybil-resistance
mechanism. A Sybil attack occurs when a single participant creates many apparently independent
identities and presents them as separate users \cite{douceur2002sybil}. 
In the context of this protocol, an attacker could create many valid profiles, each with 
its own key pair, and use them to repeatedly Like posts, Follow accounts, or create large 
numbers of spam users. This differs from ActivityPub, where moderation responsibilities
are typically delegated to server administrators \cite{w3cActivityPub2018}. 
The implemented protocol avoids requiringa server to approve actions before they are accepted,
reducing dependence on centralised moderation. However, this also means that clients 
lack a shared protocol-level mechanism for filtering abusive, spam or coordinated accounts. 

Potential approaches include signed blocklists, trust lists or reputation systems that 
allow communities to share moderation decisions without changing the underlying 
signature-verification model. Such approaches would not necessarily prevent Sybil
attacks entirely, but could help reduce the visibility and impact of malicious accounts. 
These possibilities are discussed further in the Future Work section.

These limits distinguish what the current evaluation establishes from
security features that remain outside the prototype. The tested scenarios
show detection of the specific tampering, impersonation and repository
replacement attempts described above. Confidentiality, moderation and
resistance to coordinated account creation remain open areas for future
work.

\subsection{Limitations}

The evaluation in Chapter~\ref{sec:results} demonstrates the prototype's
main workflows, but it does not establish how the system would perform or be
used at a larger scale or by a broad range of users. The limitations below
concern the current implementation and the scope of that evaluation. Some,
such as directory trust and key recovery, relate to the security assumptions
discussed earlier in this chapter; others concern usability, the available
social interactions, publication and growth over time.

\subsubsection{Command-line workflow and account setup}\label{sec:limitation-command-line}
Although published profiles and feeds can be viewed through a standard web browser, 
creating social actions and synchronising data currently requires use of the command-line client. 
The discussion in this section therefore concerns the workflow needed to create, publish and replicate 
content, rather than the experience of viewing published content.

The client provides commands for creating identities, posting, replying,
liking, following, replicating and publishing. This gives users a layer
above issuing Git commands directly: unlike Gitweets, users do not need to
run commands such as \texttt{git push} for each social action
\cite{gitweets2026}. Social4Git also uses a command-line workflow, for
example, \texttt{social4git sync} to fetch followed users' posts
\cite{social4git2026}. GitSocial likewise requires users to run its client
to create actions and fetch followed repositories \cite{gitsocial2026}.
The prototype therefore abstracts some Git operations, but still requires
users to interact through a command-line client. To get started, users must obtain the
project, ensure Python and Git are installed, install the client's Python
dependencies, navigate to the \texttt{client} directory, and run
\texttt{python app.py}. This prints the available commands and example
invocations; users then enter commands such as
\texttt{python app.py post "Hello world"} themselves. The setup also
requires configuring the client to select the local identity and repository
created for the account. The command list and examples make the syntax
visible, and Chapter~\ref{sec:results} shows
the commands used to create and publish an account. However, users still
need to know which command to run, supply its arguments correctly, and
interpret the terminal output. Error messages and the lack of immediate,
visual feedback may be unfamiliar or inconvenient for people accustomed to
social-networking applications. The prototype does not provide a graphical
interface where users can create actions and see their effect in one place.
A future interface could guide these operations while preserving the
client-side signing and Git-based exchange described in
Chapter~\ref{sec:system-design}.

Keeping profile creation and publication separate allows an identity to be
created before a hosting repository is available, but it adds an extra step
for users who wish to publish immediately. Publication also depends on a
repository URL provided by the user. When another user follows the account,
the client compares the supplied URL with the directory entry and then
checks that the signed profile contains the same URL. These checks are shown
in Listings~\ref{lst:follow-directory} and~\ref{lst:follow-profile} in
Chapter~\ref{sec:implementation}. They confirm that the recorded URLs
agree, but do not determine whether the user originally chose the intended
repository or guarantee that it will remain available over time. An
incorrect URL may therefore cause content to be published to the wrong
repository. A URL error encountered during development demonstrated this
practical risk.
A guided setup could ask users which Git host they want to use, then help
them select an existing repository or create one and fill in its URL
automatically. Where supported, the client could use that provider's API to
create the repository after the user authorises it. This would reduce the
need to copy and paste a long URL and lower the chance of a typing mistake.
The client should still allow users to provide a repository URL from another
Git host; otherwise, convenience through one provider's API could weaken the
hosting flexibility that the protocol is intended to provide.

\subsubsection{Feed publication}\label{sec:limitation-feed-publication}

The static site is not updated automatically when a user creates or
replicates actions. As shown in the publication workflow in
Chapter~\ref{sec:results}, the user must run \texttt{publish-site} to
regenerate \texttt{feed.json} and push the resulting files to the
\texttt{gh-pages} branch. The page reflects the local actions available at
the time of publication; later actions appear after the site is published
again. This is a consequence of the prototype's static-site workflow rather
than a failure of action replication, which remains a separate process.
This is related to the command-line limitation above, but concerns a
separate step: even if users create actions through a graphical interface,
the published static page remains unchanged until its files are generated
and pushed again. A future interface could show that the site is out of date
and offer a clear way to publish the current feed. This would explain the
snapshot behaviour and simplify the extra publishing step; it would not
make a static page update itself whenever a new action is created.

\subsubsection{Social vocabulary and presentation}\label{sec:limitation-social-vocabulary}

The preceding subsection, \emph{Social Vocabulary and Data Representation},
explains why the prototype stores actions as structured JSON objects. The
current vocabulary contains four types: Posts, Replies, Likes and Follows.
This extends beyond Gitweets' post-focused workflow, which has no built-in
Follow action \cite{gitweets2026}. Social4Git's reviewed design does not include
structured replies, likes or reposts \cite{social4git2026}.
However, the prototype does not yet provide actions for editing or deleting
a published post, or for formally
resharing another user's post. ActivityPub includes Update and Delete
activities \cite{w3cActivityPub2018}. GitSocial's GitMsg format includes
reposts and quotes and can represent nested threads \cite{gitsocial2026}.

These systems represent and distribute interactions differently from this
prototype. ActivityPub sends activities through server-managed inboxes and
outboxes \cite{w3cActivityPub2018}. GitSocial records the action type in Git
commit messages \cite{gitsocial2026}. The prototype instead stores each
signed action as a JSON object in the user's repository. This structured
format gives a basis for adding explicit edit, delete, repost or quote
action types without putting social metadata in commit messages. Each new
type would still need rules for what it
references and how clients verify, replicate and display it. At present, users can mimic 
reposting or quoting by copying another person's text into a new post. 
However, this creates a separate post attributed to the copying user and does not preserve a formal, 
verifiable reference to the original signed post. Consequently, clients cannot recognise the relationship 
between the two posts. Dedicated repost and quote action types would provide this linkage and allow 
clients to display these interactions appropriately. These action types remain future work.

Replies can reference other Replies as well as Posts in the action data, so
the protocol can represent a reply to a reply. However, the current feed
renders replies directly under posts and does not display a fully nested
conversation. The action data can therefore express a relationship that the
current website does not yet present. Supporting nested conversations in
the feed remains future work.

The prototype supports text posts but does not implement image or video
attachments. ActivityPub represents posts as structured objects that can
include attachments, such as links to media files \cite{w3cActivityPub2018}.
The AT Protocol repository model also stores media blobs alongside social
records \cite{atproto2026}. By comparison, the reviewed GitSocial design
identifies storing post text in commit messages as a limitation for media
support \cite{gitsocial2026}. For this Git-based prototype, adding media would
require deciding where files are stored and how they remain available when
repositories are cloned and replicated. Git Large File Storage (Git LFS)
could be investigated as a way to manage large files alongside Git
repositories, but its hosting and availability requirements would need to
be considered.
This project did not implement or evaluate media support, so the suitability
of Git LFS remains future work.

\subsubsection{Repository growth and evaluation scale}\label{sec:limitation-repository-growth}

Git is used for large, long-lived software projects, including the Linux
kernel \cite{zhou2017scalability}; this does not by itself establish how
well this particular social-data workflow will scale. During replication,
the current implementation scans candidate action files in each configured
remote and verifies them before accepting them locally. Previously accepted
actions may therefore be checked again on later runs. As repositories
accumulate more actions, repeated scanning and signature verification may
increase processing time. No performance issue was observed in the
workflows evaluated for this thesis, but those demonstrations do not
measure performance with large repositories or many users.

Other surveyed systems handle some of this work through server-side
components. ActivityPub servers route activities through inboxes and
outboxes \cite{w3cActivityPub2018}. The AT Protocol's AppView indexes
repository updates for feed generation \cite{atproto2026}. In the
Git-as-Transport design, servers fetch repositories and scan commits to
import posts \cite{gittransport2026}. These designs can provide organised
or processed views of activity, but depend on servers to perform that work. A future
design for this protocol could investigate ways for clients to avoid
rechecking unchanged actions, or use indexes that can themselves be
verified, while retaining local signature checks. No such mechanism was
implemented or evaluated here. Performance testing with larger repositories
and more users, and the design of a suitable way to reduce repeated work,
remain future work.

The evaluation also did not examine concurrent action creation from
separate repository copies. Action identifiers use a sequential counter
based on the action files present in the local repository. If the same
identity creates actions concurrently from separate copies, those copies
may select the same next identifier. This remains a limitation of the
current action-creation workflow.

\subsection{Future Work}

The directions below build on the limitations discussed above. They describe
what could be investigated next; they are not features of the current
prototype.

The command-line and publication limitations in
Sections~\ref{sec:limitation-command-line} and
\ref{sec:limitation-feed-publication} point to one area for future work:
making the client easier to use. A simple interface could guide users
through account setup, posting, following, replication and publishing
their feed. The interface could help create a repository through a Git host
where this is available, or let users enter a repository URL from another
host. Testing this workflow with users would show whether it makes setup and
publication clearer. It would also be useful to explore ways to find
accounts without relying on the same shared directory, while still checking
that a handle points to the right key and repository.

The repository-growth limitation in
Section~\ref{sec:limitation-repository-growth} also points to work on
replication as repositories grow. The action identifiers discussed above
could be changed so that separate copies of one account can create actions
without choosing the same identifier.
This should be tested by creating actions in parallel and then merging the
copies. An index, or list of action IDs, could also tell a client which
action files to check, rather than making it scan every file each time.
Git-as-Transport uses indexes
\cite{gittransport2026}; ActivityPub organises activity in server-managed
inboxes and outboxes \cite{w3cActivityPub2018}. A Git-based design could
learn from these approaches while the client still checks each action's
signature. Since Git fetch still transfers repository data, an index
may reduce the client's repeated work without reducing the amount of data
downloaded. The workflow in Figure~\ref{fig:future-feed-index} shows one
possibility. Testing with larger repositories and more users would show
whether it improves processing time.

\begin{figure}[htbp]
    \centering
    \fbox{\parbox[c][3.5cm][c]{0.9\textwidth}{\centering
    Feed-indexing workflow figure to be added.}}
    \caption{Possible use of an index while retaining client-side action verification}
    \label{fig:future-feed-index}
\end{figure}

The action and feed limitations in
Section~\ref{sec:limitation-social-vocabulary} suggest another direction:
supporting more kinds of interaction. ActivityPub
includes Update and Delete activities \cite{w3cActivityPub2018}, while
GitSocial has repost and quote actions \cite{gitsocial2026}. Future work
could define similar signed actions that link back to the original post.
It could also improve how the feed displays replies to replies, which the
current action data can already represent. Media such as images and video
would need a clear way to store files, link them to posts and show when a
file is no longer available. Git LFS could be tested as one option for
larger files, including whether it works with different Git hosts.

The security limitations in Section~\ref{sec:limitation-security} point to
further work on account recovery, private posts and moderation. A recovery
process should let someone who loses a key prove that they can use a new
one without allowing another person to take
over their handle. Private posts could use one shared encryption key (a
session key) for a group of related posts, as explored in the archived
design notes. sAT provides an example of encrypted content on static hosting
\cite{sat2026}. The current identity
keys sign actions but do not encrypt them, so a design would need to explain
how recipients get encryption keys, what happens when access is removed,
and what information remains visible in the repository. The
workflow in Figure~\ref{fig:future-confidentiality} illustrates this idea.
Moderation could let users keep their own block choices private and choose
whether to use signed recommendations from others. These choices could
change what appears in a feed, but would not stop unwanted actions from
being replicated or prevent users from creating many accounts.

\begin{figure}[htbp]
    \centering
    \fbox{\parbox[c][3.5cm][c]{0.9\textwidth}{\centering
    Confidentiality workflow figure to be added.}}
    \caption{Illustrative workflow for optional encrypted posts and recipient key distribution}
    \label{fig:future-confidentiality}
\end{figure}

\subsection{Conclusion} Revisit: \begin{itemize} \item Research question. \item Objectives from Chapter 1. \end{itemize} Summarise: \begin{itemize} \item Main protocol design contributions. \item Prototype implementation contributions. \item Evaluation findings. \item Security findings. \end{itemize} Conclude with: \begin{itemize} \item Overall effectiveness of the Git-native approach. \item Key limitations. \item Potential future impact and directions. \end{itemize}
